\begin{frame}{수학과 Python 사전지식 리뷰 --- 세 영역}
  이번 학기는 다음 세 영역의 수학을 \textbf{계속 재사용}한다 ---
  낯선 표기가 나올 때마다 이 절로 돌아와 확인할 것:
  \begin{itemize}
    \item \kb{선형대수}: 벡터의 \kb{내적}
      $w^\top x = \sum_j w_j x_j$, \textbf{행렬-벡터 곱}, 그리고
      \kb{그래디언트}(다변수 함수를 각 변수로 편미분해서 만든
      벡터, $\nabla_w J$). 고유벡터/고유값은 Week 14(PCA)에서
      다시.
    \item \kb{미적분}: \kb{연쇄법칙 (chain rule)} --- 합성함수
      $f(g(x))$ 의 미분이 $f'(g(x))\,g'(x)$ 라는 것.
      \textbf{Week 9(역전파)의 핵심 도구}가 바로 이것.
    \item \kb{확률}: \textbf{조건부 확률}과 \kb{베이즈 정리}
      \[P(y \mid x) = \frac{P(x \mid y)\,P(y)}{P(x)}\]
      Week 3에서 본격적으로 사용.
  \end{itemize}
\end{frame}
